% Translator's report; none of the proofs below is presented as recovered author text.
\begingroup
\section*{M19-8E：8.18之后至第8节末的法英回校}
\markboth{第十九版：第8节末回校}{第十九版：第8节末回校}
\addcontentsline{toc}{section}{M19-8E：8.18之后至第8节末的法英回校}

\noindent\textbf{范围及证据。}
本次替换第十八版实际源码中 DOT\_IMPORT 11、12 两段。
起点为法文PDF物理50页最后一段“Il reste à expliciter”，经过51页，止于52页末
关于把讨论移入下一节的作者安排。对读英文PDF物理34页末、35页及36页的第9节标题以前。
法53页已开始第9节；查看53、54页只用于确定边界及核对后续是否另列假设，
不是声称此报告覆盖它们的全部译校。
本次实际渲染并查看了上述法英原页，且对法51的投影定义、逆命题，法52的域与维数条件作了放大查看。
完整原PDF、旧模块和新文件的散列，连同图像文件位置，记录在同名JSON审计文件中。

\subsection*{1. 可直接据法稿恢复的英译误读}

法50明确回指8.18，不是英34的1.18。法51的参数线性空间为 $L\subset P^\vee$，
其公共线性中心是 $L^0\subset P$。法稿明确保留 $F^0\to E$，
并在 $f^{-1}(P\setminus L^0)$ 中保留上标0。关于一般截面稠密性的句子也有截面下标；
英35漏掉这些下标，不能继续把它当作法稿的字迹缺失。

法52的函数域子扩张是小写的 $k(a_1,\ldots,a_{p-1})$；在中心上消失的形式也明确指
$L^0$。英35的 $K(a_1,\ldots,a_{p-1})$ 及 $L$ 因此已获底本证据订正。
中心维数是 $q$，不是英36的 $g$；末句的主体为 $X$ 在 $x$ 处非锥形，不是点 $x$ 在自身处非锥形。
法51直线 $D$ 旁的底域和法52线性子空间条件的底域都是 $\overline k$；英35漏了横线。

这些读法恢复不等于论证成立。特别是一般截面的稠密性、$q+1$ 与 $q+2$ 的区别，
以及到像上的投影和复合态射 $g$ 的纤维维数，须分别作数学核查。

\subsection*{2. 固定分支反例：仅有投影像二维还不够}

以下是译校反例，取任意代数闭域 $k$。令
\[
 X=\mathbf P^2_k,\qquad
 f([u:v:w])=[u^2:uv:uw:v^2:vw:w^2]\in\mathbf P^5_k,
\]
即二次 Veronese 闭浸入。以 $z_0,\ldots,z_5$ 记靶空间坐标，取
\[
 L^0=V(z_0,z_1,z_2)\simeq\mathbf P^2_k.
\]
其逆像的底层集合是固定直线 $B=V(u)$，而 $X'=D_+(u)\simeq\mathbf A^2_k$。
包含 $L^0$ 的超平面由 $[a:b:c]$ 参数化，其拉回的方程为
\[
 a u^2+buv+cuw=u(au+bv+cw)=0.
\]
在参数泛点上，$au+bv+cw$ 不与 $u$ 成比例；截面因而是两条不同直线的并。
在任意特征下这个因式分解仍成立，所以泛截面可约。
但在 $u\ne0$ 上，投影复合态射为
\[
 g([u:v:w])=[u^2:uv:uw]=[u:v:w],
\]
其像为 $Q=\mathbf P^2$ 的稠密仿射开集，故 $\dim g(X')=2$。
$X'_\xi$ 只是移动直线 $V(au+bv+cw)$ 的非空开部，在整个两直线并中不稠密。

此例否定的是法51把任意中心推广到一般线性系时遗漏基点条件的断言。
它甚至不依赖于收缩态射，因为 $f$ 是闭浸入。
原稿对经过一个点或在一个光滑点相切的特殊线性系的最终适用范围，仍须分别检查基点集；
本例不代替这些特殊情形的分析。

\subsection*{3. 正文所补稠密性条件的证明}

以下为译校补证。设 $X$ 是 $k$ 上有限型的几何不可约概形，$d=\dim X$，
$f(X)\not\subset L^0$。设 $B=f^{-1}(L^0)$ 满足 $\dim B\le d-2$。
在考察几何不可约性时，可以在几何基变换之后取约化结构；所得底层空间不变，
而约化的 $X$ 是整概形。

参数泛点所给的线性组合不在 $X$ 上恒为零。因此，在 $X$ 的每个平凡化开集上，
截面由一个非零函数定义。非空截面的每个不可约分支均余维1：主理想定理给出高度至多1，
整性排除高度0；有限型域代数的维数公式再给出分支维数 $d-1$。
这一结论在扩张到底域代数闭包后同样成立。
由于 $B$ 的维数至多 $d-2$，没有任何这样的分支包含在 $B$ 中。
所以 $X'_\xi=Y_\xi\setminus B_{k(\xi)}$ 在 $Y_\xi$ 中稠密，且几何基变换后仍稠密。
非空稠密开集与全空间的不可约性等价，因而二者的几何不可约性等价。

若 $\dim g(X')\ge2$，3.1应用于 $g$ 后使 $X'_\xi$ 非空且几何不可约，
上述稠密性才把结论传给整个 $Y_\xi$。如果 $\dim g(X')=0$，
一般超平面避开它的像，$X'_\xi$ 为空；这也解释了为什么原文单独的
$f(X)\not\subset L^0$ 不能保证所要求的维数。

一个可供交叉核对的独立出处是
\href{https://stacks.math.columbia.edu/tag/0G4C}{Stacks Project，37.32，Bertini theorems}。
其中引理37.32.3的条件(4)明确要求共同基点集各不可约分支的余维至少为2。
这里引用这一点以核对缺失条件的性质；上面的具体反例和稠密性证明为本次译校自行给出，
不冒称是手稿漏失的原证明。

\subsection*{4. 对偶约定与“多两维”}

本篇采用 $\mathbf P(E)$ 表示 $E$ 的一维商的约定。令 $F^0\subset E$ 为定义中心的线性形式空间，
则
\[
 L^0=\mathbf P(E/F^0),\qquad Q=\mathbf P(F^0),\qquad
 L=\mathbf P((F^0)^\vee)=Q^\vee.
\]
将 $E$ 的一个一维商限制到 $F^0$，恰在该限制非零时得到 $u$；
限制为零的点恰是 $L^0$。这些公式直接说明为什么超平面参数空间是 $L$，而投影靶是 $Q$。
法稿首次写 $Q$ 的字母经过覆打，不把这种线性代数解释冒称为对覆打笔画的完全辨读。

写 $N=\dim P$，$q=\dim L^0$，则 $\dim F^0=N-q$。
若 $D\subset Q$ 是一条直线，可写
\[
 D=\mathbf P(F^0/K),\qquad \dim(F^0/K)=2.
\]
所以 $u^{-1}(D)$ 的射影闭包为 $\mathbf P(E/K)$，其维数是
\[
 \dim E-\dim K-1
 =N+1-\bigl((N-q)-2\bigr)-1=q+2.
\]
它包含中心 $\mathbf P(E/F^0)$，而且限制投影的像正是 $D$。
故当投影像为一维时，“该像包含于一条直线”对应于原像位于一个
包含 $L^0$、比 $L^0$ 多\emph{两}维的线性子空间中。

法52的可辨原文是“contenant $L^0$ et de dimension une de plus”，即“包含 $L^0$，维数多一维”；
英35--36则写为“dimension at most one”。前者少了一维，后者还改变了比较对象。
取 $f=\mathrm{id}_{\mathbf P^2}$、$L^0$ 为一点时，$g$ 的像是 $\mathbf P^1$，
而含该点及整个 $f(X)$ 的最小线性空间是 $\mathbf P^2$。这是无需任何深层定理的直接反证。

\subsection*{5. 作者拟议的3.1逆命题：说明几何基变换后的版本}

本小节为译校补证，不是法稿原有的证明。
几何不可约性允许先扩张到 $\overline k$，再将 $X$ 及像闭包取约化结构。
以 $Z$ 表示这个不可约像闭包，泛超平面截面仍记作 $Y$。
若 $\dim Z=0$，泛超平面避开 $Z$，故 $Y$ 为空，并非几何不可约。
若 $\dim Z\ge2$，结论就是3.1。

剩下 $Z$ 是整射影曲线的情形。令 $e=\deg Z$。
一般超平面避开 $Z$ 的有限奇异点集，也避开 $Z$ 在其光滑点处的切超平面族，
因而在代数闭域上同 $Z$ 相交于 $e$ 个不同的点。
后一切超平面族的维数至多为对偶射影空间维数减一：在每个光滑点，
须包含其射影切线的超平面族余维2，而点本身只提供一维参数。
若靶为 $\mathbf P^1$，该族为空，结论同样成立。
此外，$f(X)$ 是 $Z$ 的稠密可构造子集，因而包含稠密开子集；
一般超平面还可以避开其有限补集。所以这 $e$ 个点都有非空逆像。

当 $e>1$，$Y$ 在几何参数泛点上分成这 $e$ 个不同点的非空逆像；
它们彼此不交，且在 $Y$ 中均开闭，因此 $Y$ 不可约性不成立。
当 $e=1$，整射影曲线 $Z$ 就是一条直线 $D$。
一般超平面与 $D$ 相交于唯一一点。描述“点在超平面上”的关联空间，在 $D$ 上是射影丛；
在对偶空间上则双有理于一般超平面的参数空间。
因此，它的函数域在 $k(D)$ 上是纯超越扩张，而泛超平面截面由 $X\to D$ 的泛纤维作该扩域取得。
几何不可约性在扩域后保持并可下降，故两者几何不可约性等价。

这证明了正文采用的几何版本：像维数至少2，或者像为直线且泛纤维几何不可约。
法51第一分支的符号有改动且形似“$\ne2$”；英35明确印成“$=2$”。
不能声称从原字迹直接辨出了“$\ge2$”，但 $f=\mathrm{id}_{\mathbf P^3}$
已排除“$=2$”作为正确数学条件。法稿的直线标有 $\overline k$，正文因而统一在几何基变换后陈述，
并不擅自加强为直线必须在原底域 $k$ 上定义。

\subsection*{6. 正确的纤维维数估计及其界限}

投影 $u:P\setminus L^0\to Q$ 的每个纤维同构于 $\mathbf A^{q+1}$。
令 $W=\overline{f(X)}$，$W'=W\setminus L^0$。
限制态射 $W'\to Q$ 的纤维为这些仿射空间的闭子概形，维数至多 $q+1$。
又因为 $f(X')$ 在 $W'$ 中稠密，它的投影像闭包同 $u(W')$ 的像闭包相同。
维数公式于是给出
\[
 \dim f(X)=\dim W'\leq\dim g(X')+q+1.
\]
这正是原稿所需的不等式。

这里不能将 $W'\to Q$ 换成 $g:X'\to Q$ 来断言纤维上界不变。
取 $X=\mathbf P^N\times\mathbf P^m$、$m>0$，$f$ 为第一投影；
若中心为 $\mathbf P^q$，则 $g$ 的纤维是
$\mathbf A^{q+1}\times\mathbf P^m$，维数为 $q+1+m$。
法52的手写补句和英36都写成“$g$ 的纤维维数至多 $q+1$”，这是两本共有的错误，
故正文订正论证对象，而不改变可保留的维数不等式。

对 $q=0$，$\dim f(X)\ge3$ 只保证去掉 $f^{-1}(L^0)$ 后的泛截面几何不可约。
若要断言整个泛截面不可约，仍须控制固定分支。
例如 $\operatorname{Bl}_p\mathbf P^3\to\mathbf P^3$ 中，经过 $p$ 的任一超平面总变换为其严格变换与例外除子之和；
两者是不同的非空分支，所以8.17(a)在其放宽后的宽泛假设下不能断言坏参数不包含超平面。
这不反驳8.18另外规定的非分歧性范围，也不证明本版加强后的8.18已从本段推出。
最后“在 $x$ 处不是锥形”的句子未给出定义与证明，保留为作者提示。

\subsection*{7. 删改、未辨字与历史脚注}

本次保留的可辨删改包括：法50将整性及代数闭底域改为几何不可约性；
法51将 \emph{simple} 改为 \emph{lisse}，将截面的几何整条件改为几何不可约，
将 $Q$ 的旧商式改为 $\mathbf P(F^0)$，并改写导出投影的线性映射；
法52划去形容“条件”的 \emph{fondamentale}，且手写插入泛纤维条件与纤维维数补句。
其余被连续叉号盖住的短词或旧半句不具备可靠逐字读法，图像中照存；不根据后文重造被删句。
“cunutesque”的意义、法51首次 $Q$ 字母覆打和逆命题维数符号的最终字形仍未可靠解决。

旧译脚注全部随原第十八版源码保留；下列留存只为方便本轮对照，所有“底本未明”均是旧译当时的判断，
不覆盖本次法稿回校已解决的读法。当前采用的结论以前面正文脚注和本报告为准。
\begin{enumerate}
\item 底本印作“1.18”，但紧接8.18并明确指向该命题的(a)、(b)，此处订为8.18。

\item 底本此处写“a linear subariety $L$ of $P$”，随后又用 $L$ 表示被这些超平面包含的线性子簇，并在后文写作 $L^0$。为区分两个不同空间，正文将参数线性系称为 $L$，其基线性子空间称为 $L^0\subset P$；此举不补定后面关于 $F^0$ 的未明约定。

\item 底本写成“$f(X)\notin L-0$”，下一句立即给出上述逆像稠密条件；这里订正为像不包含于中心 $L^0$，不是把像当成一个点。

\item 这一句的射影对偶约定尚不清楚：原稿同时把 $Q$ 称为超平面参数空间，又在下面使用 $L\simeq Q^\vee$，二者不能不加解释地混同。$F^-$ 的上标也与前面的 $F^0$ 不同。此处译校保留这些写法并标待核，不据此猜补向量空间的商、对偶或映射。

\item 底本第二条显示公式写 $f^{-1}(P-L)$，但前面的 $X'$ 与 $u$ 的定义域均为中心 $L^0$ 的补集；此处统一为 $L^0$。$f'$ 指 $f$ 在 $X'$ 上的限制。

\item 底本在这句话中未写 $\xi$ 下标，但上句正在讨论截面族 $X'_\xi$。可能本意是相应截面中的稠密性，而不是重复全空间 $X'\subset X$；因原稿未明写，正文引述其现存记号，不擅自补成对所有截面的断言。

\item 底本确实为“$\dim f(X)=2$”，不是“$\ge2$”；它与紧前一句所用的充分条件存在明显张力。这里引述作者拟议的逆命题并标待核，不能把“当且仅当”当成已核实的定理，也不在缺少完整假设和证明时擅自补成一个修订逆命题。

\item 这句话依赖上面尚待核的逆命题；忠实译出不表示已验证该充要条件。

\item 这里“包含 $X$”沿用底本措辞；在给定态射 $f:X\to P$ 的设置下，应按其像 $f(X)$ 理解，而不是假定 $f$ 已是嵌入。

\item 底本在这句写“vanishing on $L$”，但本段中心为 $L^0$，且紧接着指定 $p=\operatorname{codim}(L^0,P)$；这里统一中心记号为 $L^0$，区别于前段的超平面参数线性系 $L$。在所选仿射开集上，线性形式经平凡化成为函数，再拉回到函数域。

\item 底本写为“$K(a_1,\ldots,a_{p-1})\subset K$”。由于 $a_i\in K$，该域就是 $K$，一般不能给出像的维数。正文订为 $k(a_1,\ldots,a_{p-1})$：在投影的首坐标非零的仿射图中，像的坐标环由这些坐标比生成，其分式域即此子域，故像闭包的维数等于它对 $k$ 的超越次数。这一订正仅针对本句，并不恢复前批关于 $Q$ 的全部对偶约定。

\item “dimension at most one”是底本跨第35、36页的现存文字，不能直接当成已核实判据。若 $\dim L^0>1$，这样的线性子簇根本不能包含 $L^0$；而从中心 $L^0$ 投影后像落在直线的条件，通常对应于其逆像落在更高维的线性子空间。这里只明确引述原稿要求，未猜补可能漏失的相对维数或改写为未经证实的充要条件。

\item 底本先写“dimension $g$”，下一行及后续却均用 $q$；此处将维数变量统一为 $q$，不与态射 $g$ 混同。这里采用前文由中心 $L^0$ 给出的线性投影。

\item 底本写“so that those of $g$ are $\dim\le q+1$”。一般 $g=u\circ f'$ 的纤维还包含 $f$ 自身的纤维，不能仅由 $u$ 的纤维界直接得到这一说法。正文改用投影在像 $f(X)\setminus L^0$ 上的限制：其纤维是 $u$ 的纤维中的子集，维数界及随后关于像维数的不等式成立。若另知 $f$ 拟有限，原稿对 $g$ 的纤维界亦可成立；此处不擅自新增该假设。像的维数按其闭包的维数理解，并沿用 $X'=f^{-1}(P-L^0)$ 稠密的前提。

\item 底本为“provided $x$ is not ‘conical at $x$ relative to $f$’”。首个小写 $x$ 可能本应为 $X$，也可能是在描述点 $x$ 为相对于 $f$ 的锥点；对象和定义均未明，正文保留现存措辞，不确定恢复。原稿未在本段给出精确定义或完整论证，故这句不作为已经独立验证的判据。

\end{enumerate}
\par
\endgroup
